        % ============================================================
%  Plantilla LaTeX para Trabajos de Título Escuela de Ingeniería
%  Adaptada desde la plantilla oficial en Word
% ============================================================
\documentclass[12pt,letterpaper]{report}

% ---------- Idioma y codificación ----------
\usepackage[utf8]{inputenc}
\usepackage[T1]{fontenc}
\usepackage[spanish]{babel}

% ---------- Márgenes  ----------
\usepackage[top=2.54cm,bottom=2.54cm,left=2.54cm,right=2.54cm]{geometry}

% ---------- Tipografía, Escritura de Fórmulas y Códigos de Fuente ----------
\usepackage[defaultsans]{opensans} % Tipo de letra del documento (solo reemplaza sans-serif, más liviano que 'default')
\usepackage{sansmathaccent}
\usepackage{amsmath}
\usepackage{amssymb}
\usepackage{amsthm}

% ---------- Código fuente (listings en vez de minted: sin shell-escape, compila mucho más rápido) ----------
\usepackage{listings}
\usepackage{xcolor}

\definecolor{uohblue}{HTML}{1F7AC2}
\definecolor{uohgray}{HTML}{595959}
\definecolor{codebg}{HTML}{F7F7F7}
\definecolor{codegreen}{HTML}{2E7D32}

\lstdefinestyle{pythonstyle}{
    language=Python,
    backgroundcolor=\color{codebg},
    basicstyle=\ttfamily\small,
    keywordstyle=\color{uohblue}\bfseries,
    commentstyle=\color{uohgray}\itshape,
    stringstyle=\color{codegreen},
    numbers=left,
    numberstyle=\tiny\color{gray},
    numbersep=8pt,
    breaklines=true,
    breakatwhitespace=false,
    showstringspaces=false,
    frame=single,
    rulecolor=\color{uohgray},
    tabsize=4,
    captionpos=b
}
\lstset{style=pythonstyle}

% ---------- Paquetes para tablas ----------
\usepackage{longtable}
\usepackage{booktabs}
\usepackage{multirow}

% ---------- Interlineado y espaciado de párrafo ----------
\usepackage{setspace}
\onehalfspacing % interlineado 1.5, similar a la lectura cómoda del original
\usepackage{indentfirst} % sangría en la primera línea de cada párrafo (como el original)
\setlength{\parindent}{1.25cm}
\setlength{\parskip}{0pt}

% ---------- Imágenes ----------
\usepackage{graphicx}
\graphicspath{{img/}}

% ---------- Hipervínculos ----------
\usepackage{hyperref}
\hypersetup{
    colorlinks=true,
    linkcolor=black,
    urlcolor=uohblue,
    citecolor=black,
    pdftitle={Evaluación y fortalecimiento de la seguridad de software y arquitectura de la plataforma Maria.ag},
    pdfauthor={Felipe Sebastián Suárez González},
    pdfsubject={Trabajo de título o grado},
    bookmarksnumbered=true
}

% ---------- Títulos de sección ----------
\usepackage{titlesec}
% Estilo de los capítulos/secciones: centrado, en negrita, sin numeración
% (el Word original no numera las secciones del cuerpo)
\titleformat{\chapter}[block]
  {\centering\bfseries\large}
  {}{0pt}{}
\titlespacing{\chapter}{0pt}{0pt}{24pt}

% ---------- Índice con puntos guía (estilo Word) ----------
\usepackage{tocloft}
\renewcommand{\cftchapfont}{\bfseries}
\renewcommand{\cftchappagefont}{\bfseries}
\renewcommand{\cftchapleader}{\cftdotfill{\cftdotsep}}
\setlength{\cftbeforechapskip}{6pt}
\renewcommand{\cfttoctitlefont}{\hfill\bfseries\large}
\renewcommand{\cftaftertoctitle}{\hfill}
\usepackage{fancyhdr}
\pagestyle{fancy}
\fancyhf{}
\fancyfoot[C]{\thepage}
\renewcommand{\headrulewidth}{0pt}

% ---------- Citas y referencias bibliográficas (formato APA) ----------
\usepackage[backend=biber,style=ieee]{biblatex}
\usepackage{csquotes}
\addbibresource{referencias2.bib}

% ---------- Listas (para hipótesis, objetivos, preguntas con viñetas) ----------
\usepackage{enumitem}


% ============================================================
\begin{document}
% ------------------------------------------------------------
% PORTADA
% ------------------------------------------------------------
\begin{titlepage}
    \centering
    \vspace*{1cm}

    \includegraphics[width=7cm]{logo_uoh.png}\\[0.3cm]
    {\large Escuela de Ingeniería\\
    Ingeniería Civil en Computación}\\[2.5cm]

    \begin{minipage}{0.8\textwidth}
        \centering
        {\bfseries\Large Evaluación y fortalecimiento de la seguridad de software y arquitectura de la plataforma Maria.ag}
    \end{minipage}\\[2cm]

    {\large Felipe Sebastián Suárez González}\\[0.3cm]
    {\large Profesor guía: Edgardo Ortiz}\\[2cm]

    {Memoria para optar al título de ingeniero civil en computación}\\

    \vfill

    {Rancagua, Chile\\
    2026}

    \vspace*{1cm}
\end{titlepage}
% ------------------------------------------------------------
% DECLARACIÓN USO IA
% ------------------------------------------------------------
\chapter*{Declaración de uso de Inteligencia Artificial}

\textit{[Esta sección debe completarse indicando si se utilizó inteligencia artificial generativa en la elaboración del trabajo de título. En caso afirmativo, siga el formato sugerido a continuación; en caso negativo, indique explícitamente que no se utilizaron herramientas de IA generativa.]}

\vspace{0.3cm}

\textbf{Formato sugerido (si hubo uso de IA):}

En la elaboración de este trabajo de título se utilizó \textit{[nombre de la herramienta, ej. Claude, ChatGPT, Copilot]}, desarrollado por \textit{[nombre del editor/desarrollador]}, como apoyo en las siguientes tareas: \textit{[ej. revisión gramatical y de estilo del texto; generación de código de ejemplo en la sección de Anexos; asistencia en la estructuración del marco metodológico]}. Su uso se concentró específicamente en \textit{[capítulos o secciones, ej. el Capítulo de Metodología y el Anexo B]}.

Todo el contenido generado o asistido por inteligencia artificial fue revisado, verificado y validado por el/la autor(a), quien asume la responsabilidad íntegra sobre la exactitud, originalidad y pertinencia del contenido final presentado en este documento.
% ------------------------------------------------------------
% DEDICATORIA (opcional)
% ------------------------------------------------------------
\chapter*{Dedicatoria (opcional)}
Introducir opcionalmente una frase en la que se dedique el trabajo de título o grado a una o más personas particulares. Generalmente, las dedicatorias incluyen el tipo de vínculo con la(s) persona(s) a quien(es) se dedica el trabajo, su nombre de pila y la razón por la que se dedica el texto a esta persona. Por ejemplo: ``A mi hermana Patricia, por su apoyo constante''.

\newpage

% ------------------------------------------------------------
% AGRADECIMIENTOS (opcional)
% ------------------------------------------------------------
\chapter*{Agradecimientos (opcional)}
Introducir opcionalmente los agradecimientos que sean pertinentes ya sea a individuos, organizaciones, entre otros. Generalmente, los agradecimientos son presentados en primera persona, identificando el acto de habla ``agradecer'', las personas, grupos u organizaciones a las que se agradece y la razón. Por ejemplo: ``Agradezco especialmente a los(as) participantes de este estudio, quienes hicieron posible (\dots). Agradezco también a mis padres, quienes con su apoyo (\dots)''.

\newpage

% ------------------------------------------------------------
% ÍNDICE
% ------------------------------------------------------------
\renewcommand{\contentsname}{Índice}
\tableofcontents

\newpage

% ============================================================
% CUERPO DEL DOCUMENTO
% ============================================================

\chapter*{Resumen}
\addcontentsline{toc}{chapter}{Resumen}

Este trabajo evalúa la seguridad del software y de la arquitectura de la
plataforma Maria.ag de la empresa Gota SpA, implementa medidas de mitigación
sobre los hallazgos priorizados y verifica su efectividad, tomando como
referencia la Guía de Pruebas de Seguridad Web de OWASP, el OWASP Top 10 y el
sistema de puntuación CVSS.

\newpage

\chapter*{Introducción}
\addcontentsline{toc}{chapter}{Introducción}

En este capítulo se presenta el contexto del trabajo realizado en esta memoria
de título, la problemática que lo motiva, su justificación y los objetivos que
se abordan.

Gota SpA es una empresa especializada en soluciones de riego avanzadas, cuyo
objetivo es mejorar el uso del agua en la agricultura, reducir costos y
aumentar la productividad de los cultivos \autocite{gota}. Entre sus servicios
cuenta con monitoreo de equipos de riego, instalación de equipos de telemetría
y un software propio, Maria.ag, que utiliza un modelo de riego predictivo para
regar de manera óptima, realiza un balance hídrico en tiempo real y entrega
recomendaciones de cuándo y cuánto regar \autocite{gota,gota2024maria}.

Maria.ag es una aplicación web multiusuario: cada empresa cliente accede
mediante cuentas con credenciales propias y datos de campo que le son
exclusivos (suelos, sectores de riego, especies cultivadas), por lo que la
plataforma debe garantizar un aislamiento efectivo entre los distintos
clientes que comparten la misma infraestructura. Sobre la base de los datos
ingresados por el usuario y de la información recibida mediante integraciones
externas, el modelo predictivo genera recomendaciones de riego que el usuario
puede aceptar o complementar registrando riegos de forma manual. Estas
recomendaciones y registros no permanecen únicamente dentro de la aplicación:
Maria.ag se comunica mediante API con plataformas de terceros --como
Dropcontrol, Talgil y Galcon-- a las que no solo consulta información
meteorológica y de telemetría, sino a las que también envía instrucciones de
riego para su ejecución. Esta última característica traslada una parte del
riesgo de seguridad del plano exclusivamente digital al plano físico, ya que
una falla en la validación de los datos que viajan hacia estas integraciones,
o un acceso indebido a las credenciales que las habilitan, podría derivar en
la activación incorrecta de equipos de riego en terreno.

Según lo informado por la empresa, durante el desarrollo de la plataforma se
priorizó alcanzar una operatividad mínima, dejando en segundo plano aspectos
como la documentación o la realización de pruebas rigurosas de seguridad.
Esta situación expone a la plataforma a posibles vulnerabilidades cuyo
impacto puede ir desde la filtración o alteración de información de los
clientes --incluyendo el acceso de una empresa a los datos de otra-- hasta la
alteración del funcionamiento de la plataforma y, en el caso de las
integraciones con capacidad de accionamiento, del propio proceso de riego en
campo.

Frente a este escenario, este trabajo propone realizar una evaluación de
seguridad de la arquitectura y del software de Maria.ag, con especial
atención al aislamiento entre clientes, a la autenticación y gestión de
sesiones, a la validación de los datos ingresados y transmitidos, y a la
seguridad de las integraciones con plataformas externas, seguida de la
implementación de medidas de mitigación y de la verificación de su
efectividad. Para ello, se utilizará como referencia la Guía de Pruebas de
Seguridad Web de OWASP \autocite{owasp2024wstg}, el OWASP Top 10
\autocite{owasp2025top10}, el sistema de puntuación CVSS
\autocite{first2023cvss} y, dado el carácter de las integraciones externas de
la plataforma, los lineamientos del OWASP API Security Top 10.

El trabajo se justifica en la necesidad de proteger los datos de los clientes,
la continuidad operativa de la aplicación, la integridad del proceso físico
de riego que depende de ella y, según lo informado por la empresa, la
ausencia de evaluaciones de seguridad previas en la empresa.

El objetivo general y los objetivos específicos se presentan en los capítulos
siguientes.

\newpage

\chapter*{Hipótesis}
\addcontentsline{toc}{chapter}{Hipótesis}

La aplicación de un proceso sistemático de evaluación de seguridad basado en metodologías y estándares reconocidos, junto con la posterior implementación de medidas de mitigación, permite reducir el número y la severidad de las vulnerabilidades en la plataforma Maria.ag.
\newpage

\chapter*{Preguntas de investigación}
\addcontentsline{toc}{chapter}{Preguntas de investigación}

\begin{itemize}[leftmargin=1.25cm]
    \item ¿Cuáles son los principales activos, componentes y superficies de ataque de la plataforma Maria.ag?
    \item ¿Qué vulnerabilidades y riesgos presenta la plataforma según la metodología WSTG, el OWASP Top 10 y la clasificación CVSS?
    \item ¿Qué criterios permiten priorizar las vulnerabilidades para su mitigación dentro del tiempo disponible?
    \item ¿En qué medida las mitigaciones implementadas reducen el número y la severidad de los hallazgos?
\end{itemize}

\newpage

\chapter*{Objetivo general}
\addcontentsline{toc}{chapter}{Objetivo general}

El objetivo principal de este trabajo es evaluar y fortalecer la seguridad del software y de la arquitectura de la plataforma Maria.ag de la empresa Gota SpA, mediante la identificación de vulnerabilidades, la implementación de medidas de mitigación y la verificación de su efectividad, tomando como referencia la Guía de Pruebas de Seguridad Web de OWASP \autocite{owasp2024wstg}.
\newpage

\chapter*{Objetivos específicos}
\addcontentsline{toc}{chapter}{Objetivos específicos}

\begin{enumerate}[leftmargin=1.25cm]
    \item Caracterizar la arquitectura, los componentes, las dependencias, los activos y los mecanismos de seguridad de la plataforma Maria.ag, delimitando su superficie de ataque y los requerimientos críticos del sistema.
    \item Identificar y clasificar vulnerabilidades y riesgos técnicos mediante modelado de amenazas, análisis estático y dinámico, revisión de dependencias y evaluación de la arquitectura, utilizando como referencia la Guía de Pruebas de Seguridad Web de OWASP \autocite{owasp2024wstg}, el OWASP Top 10 \autocite{owasp2025top10} y el Sistema de Puntuación de Vulnerabilidades Comunes (CVSS) \autocite{first2023cvss} para la clasificación de severidad.
    \item Priorizar los hallazgos según su severidad, impacto y factibilidad de intervención, y diseñar e implementar medidas de mitigación para un conjunto seleccionado de vulnerabilidades y problemas arquitectónicos.
    \item Evaluar las mejoras implementadas mediante la repetición de las pruebas de seguridad y la comparación de indicadores antes y después de la intervención (número de hallazgos por severidad y puntajes CVSS), complementando los resultados con pruebas de regresión.
\end{enumerate}
\newpage

% ------------------------------------------------------------
% MARCO TEÓRICO
% ------------------------------------------------------------
\chapter*{Marco teórico}
\addcontentsline{toc}{chapter}{Marco teórico}

En este capítulo se presentan los conceptos, estándares y metodologías que
sustentan la evaluación de seguridad desarrollada en este trabajo. Se abordan,
en orden, los fundamentos de la seguridad de aplicaciones web, las
metodologías y marcos de referencia utilizados para guiar la evaluación, los
riesgos propios de arquitecturas multiusuario como la de Maria.ag, la
seguridad de las interfaces de programación de aplicaciones (API) dado el rol
central que estas cumplen en la plataforma, y las particularidades de
seguridad que surgen cuando un sistema de software tiene capacidad de incidir
sobre procesos físicos, como ocurre con el riego agrícola.

\section*{Seguridad de la información y de las aplicaciones web}
\addcontentsline{toc}{section}{Seguridad de la información y de las aplicaciones web}

La seguridad de la información se organiza tradicionalmente en torno a tres
propiedades conocidas como la tríada CIA: \textit{confidencialidad}, que
garantiza que la información solo sea accesible para quienes están
autorizados a conocerla; \textit{integridad}, que asegura que los datos no
sean alterados de forma no autorizada o accidental; y \textit{disponibilidad},
que asegura que los sistemas y datos estén accesibles cuando se requieran
\autocite{stallings2018}. Estas tres propiedades constituyen el criterio base
para clasificar el impacto de cualquier vulnerabilidad identificada durante la
evaluación.

En el caso particular de las aplicaciones web, la seguridad se ve condicionada
por su arquitectura cliente-servidor y por la exposición pública de sus
interfaces, lo que amplía la superficie de ataque respecto de aplicaciones de
uso exclusivamente local. Los riesgos más comunes en este tipo de sistemas
incluyen fallas en el control de acceso, en la validación de entradas, en la
gestión de sesiones y en la configuración de los componentes que conforman la
aplicación \autocite{owasp2025top10}.

\section*{Metodologías y estándares de referencia}
\addcontentsline{toc}{section}{Metodologías y estándares de referencia}

\subsection*{OWASP Top 10}
El \textit{Open Worldwide Application Security Project} (OWASP) es una
comunidad abierta dedicada a mejorar la seguridad del software. Su proyecto
OWASP Top 10 \autocite{owasp2025top10} identifica y describe las diez
categorías de riesgo más críticas en aplicaciones web, construidas a partir de
datos aportados por organizaciones de seguridad a nivel mundial. Entre estas
categorías se encuentran el control de acceso roto (\textit{broken access
control}), las fallas criptográficas, la inyección de código y las fallas de
diseño relacionadas con la lógica de negocio de la aplicación. Este trabajo
utiliza el OWASP Top 10 como taxonomía para clasificar los hallazgos
identificados.

\subsection*{OWASP Web Security Testing Guide (WSTG)}
La Guía de Pruebas de Seguridad Web de OWASP \autocite{owasp2024wstg} es un
marco metodológico que describe, de manera estructurada, los procedimientos
para evaluar la seguridad de una aplicación web a lo largo de su ciclo de
desarrollo. La guía organiza las pruebas en categorías --gestión de la
configuración, autenticación, gestión de sesiones, validación de datos,
lógica de negocio, entre otras-- y, para cada una, define objetivos, técnicas
y criterios de verificación. En este trabajo, la WSTG se utiliza como
procedimiento operativo para ejecutar la evaluación de seguridad de
Maria.ag.

\subsection*{Common Vulnerability Scoring System (CVSS)}
El Common Vulnerability Scoring System \autocite{first2023cvss} es un
estándar abierto que permite asignar una puntuación numérica a la severidad
de una vulnerabilidad, a partir de métricas base --como el vector de ataque,
la complejidad de explotación y el impacto sobre la confidencialidad,
integridad y disponibilidad-- que pueden complementarse con métricas
temporales y de entorno. La puntuación resultante, en una escala de 0 a 10,
permite comparar y priorizar hallazgos de distinta naturaleza bajo un mismo
criterio cuantitativo, lo que en este trabajo se utiliza tanto para priorizar
las vulnerabilidades a mitigar como para comparar el estado de la plataforma
antes y después de la intervención.

\subsection*{Modelado de amenazas}
El modelado de amenazas es un proceso estructurado para identificar, de forma
anticipada, las posibles formas en que un sistema puede ser atacado, a partir
de la descomposición de su arquitectura, la identificación de sus activos y
puntos de confianza, y el análisis de los flujos de datos entre sus
componentes. Un marco ampliamente utilizado con este fin es STRIDE, que
clasifica las amenazas en seis categorías: suplantación de identidad
(\textit{Spoofing}), alteración de datos (\textit{Tampering}), repudio
(\textit{Repudiation}), divulgación de información (\textit{Information
disclosure}), denegación de servicio (\textit{Denial of service}) y elevación
de privilegios (\textit{Elevation of privilege}) \autocite{shostack2014}. En
este trabajo, el modelado de amenazas se utiliza como insumo previo a las
pruebas dinámicas, para orientar el esfuerzo de evaluación hacia los
componentes y flujos de mayor riesgo.

\section*{Control de acceso en arquitecturas multiusuario}
\addcontentsline{toc}{section}{Control de acceso en arquitecturas multiusuario}

Maria.ag opera bajo un modelo en que múltiples empresas clientes comparten la
misma instancia de la aplicación e infraestructura, cada una con sus propios
datos y usuarios. Este tipo de arquitectura, conocida como
\textit{multiusuario} o \textit{multi-tenant}, requiere que el control de
acceso se aplique de forma consistente en cada operación del sistema y no
únicamente a nivel de interfaz, de modo que un usuario autenticado solo pueda
leer o modificar los recursos que le pertenecen \autocite{owasp2025top10}.

Cuando esta verificación falla, se produce lo que OWASP clasifica como
control de acceso roto, cuya manifestación más frecuente es la referencia
directa insegura a objetos (\textit{Insecure Direct Object Reference}, IDOR):
un identificador de un recurso --por ejemplo, el de un sector de riego o un
registro de riego manual-- es manipulado por el usuario para acceder a
recursos que pertenecen a otra cuenta o empresa, sin que el servidor verifique
la propiedad del recurso solicitado. Dado el modelo multiusuario de Maria.ag,
este trabajo otorga especial relevancia a la verificación del control de
acceso entre clientes y entre roles dentro de una misma cuenta.

\section*{Seguridad de interfaces de programación de aplicaciones (API)}
\addcontentsline{toc}{section}{Seguridad de interfaces de programación de aplicaciones (API)}

Una parte sustancial de la funcionalidad de Maria.ag se sustenta en el
intercambio de información con plataformas externas a través de API, tanto
para la obtención de datos meteorológicos y de telemetría como para el envío
de recomendaciones de riego hacia sistemas de terceros. Esta dependencia hace
pertinente incorporar al marco de referencia de este trabajo los riesgos
específicos de las API, los cuales no siempre son cubiertos en su totalidad
por las guías orientadas a aplicaciones web tradicionales.

El OWASP API Security Top 10 describe los riesgos más críticos en este tipo
de interfaces, entre los que se cuentan la autorización rota a nivel de
objeto, la asignación masiva de propiedades no esperadas en una petición, la
exposición excesiva de datos en las respuestas, y la ausencia de límites de
uso (\textit{rate limiting}) que permitan contener el abuso de un endpoint.
Adicionalmente, cuando una API actúa como intermediaria hacia un sistema
externo, cobra relevancia el correcto manejo de las credenciales de dicha
integración --su almacenamiento, rotación y alcance de permisos-- así como la
validación de los datos antes de reenviarlos, para evitar que un valor
manipulado por el cliente sea transmitido sin control hacia la plataforma de
destino.

\section*{Seguridad en sistemas con componentes ciberfísicos}
\addcontentsline{toc}{section}{Seguridad en sistemas con componentes ciberfísicos}

A diferencia de una aplicación web convencional, Maria.ag forma parte de un
sistema ciberfísico: sus recomendaciones y registros de riego son transmitidos
a plataformas externas que gestionan equipamiento de riego en terreno
(bombas, válvulas, controladores), por lo que una vulnerabilidad en el
software puede manifestarse como un efecto físico sobre el cultivo o la
infraestructura de riego, y no únicamente como una filtración o alteración de
datos. Esta característica es análoga a la que se observa en los sistemas de
control industrial y en el Internet de las Cosas (IoT) aplicado a la
agricultura de precisión, donde la literatura ha documentado que errores de
validación o fallas de autenticación en la capa de software pueden escalar a
consecuencias operativas y económicas sobre el proceso físico que dicho
software supervisa o controla \autocite{fathy2023secure}. En este trabajo,
esta consideración se traduce en un énfasis particular sobre la validación de
los datos que viajan hacia las integraciones con capacidad de accionamiento, y
sobre los mecanismos que evitan la duplicación o ejecución no intencionada de
una orden de riego.

\newpage
% ------------------------------------------------------------
% REFERENCIAS
% ------------------------------------------------------------
\chapter*{Referencias}
\addcontentsline{toc}{chapter}{Referencias}

\vspace{0.5cm}

\printbibliography[heading=none]

\newpage


\end{document}